% Generated draft: requires mathematical and source review.
\documentclass[11pt]{article}
\usepackage[margin=25mm]{geometry}
\usepackage{fontspec}
\setmainfont{Latin Modern Roman}
\usepackage{amsmath,amssymb,mathtools,unicode-math}
\setmathfont{Latin Modern Math}
\usepackage{enumitem,needspace,xurl,hyperref}
\setlength{\parindent}{0pt}
\setlength{\parskip}{5pt}
\setlength{\emergencystretch}{3em}
\newcommand{\N}{\mathbb{N}}
\newcommand{\Z}{\mathbb{Z}}
\newcommand{\Q}{\mathbb{Q}}
\newcommand{\R}{\mathbb{R}}
\newcommand{\C}{\mathbb{C}}
\newcommand{\F}{\mathbb{F}}
\newcommand{\A}{\mathbb{A}}
\newcommand{\PP}{\mathbb{P}}
\newcommand{\E}{\mathbb{E}}
\newcommand{\eps}{\varepsilon}
\newcommand{\dd}{\,\mathrm{d}}
\begin{document}
\section*{UnsolvedMath ID: 30004604}
\textbf{Source problem code:} catalogue entry
\subsection{Definitions and mathematical statement}
The supplied record does not define the H-constant, specify the ambient space, or state conventions for degenerate arrangements. Consequently, the following is a conditional formulation using the projective-plane, singular-locus interpretation; its equivalence to the intended source formulation requires review.

Under this interpretation, let \(C=L_1\cup\cdots\cup L_d\subset\PP^2_{\C}\) be a reduced union of \(d\geq2\) distinct projective lines. Let \(S(C)\) be the set of points incident with at least two of these lines, and let \(m_p=\#\{i:p\in L_i\}\) for \(p\in S(C)\). Define the singular-locus H-value by
\[
 H_{\mathrm{sing}}(C)=\frac{d^2-\sum_{p\in S(C)}m_p^2}{\#S(C)}.
\]
Here \(S(C)\) is nonempty under the stated conditional assumptions. An arrangement defined over \(\Q\) is interpreted as one for which each component line has a defining homogeneous linear equation with rational coefficients; its singular points are counted after extension to \(\C\).

For the supersolvable condition, associate to \(C\) the central arrangement of planes in \(\C^3\) defined by the same homogeneous linear equations. Its intersection lattice \(L\) consists of intersections of these planes, including \(\C^3\), ordered by reverse inclusion, with rank \(r(X)=\operatorname{codim}_{\C^3}X\). An element \(X\in L\) is modular if
\[
 r(X)+r(Y)=r(X\wedge Y)+r(X\vee Y)\qquad\text{for every }Y\in L.
\]
The arrangement is called supersolvable if \(L\) has a maximal chain consisting of modular elements.

Let \(\mathcal A_{\Q}\) be the class of arrangements defined over \(\Q\) under these conventions, let \(\mathcal A_{\C}\) be the class of all such complex arrangements, and let \(\mathcal A_{\mathrm{ss}}\subseteq\mathcal A_{\C}\) consist of the supersolvable ones. The conditional extremal problem is to determine
\[
 h_{\Q}=\inf_{C\in\mathcal A_{\Q}}H_{\mathrm{sing}}(C),\qquad
 h_{\C}=\inf_{C\in\mathcal A_{\C}}H_{\mathrm{sing}}(C),\qquad
 h_{\mathrm{ss}}=\inf_{C\in\mathcal A_{\mathrm{ss}}}H_{\mathrm{sing}}(C),
\]
with infima understood in the extended real numbers, and to determine whether each infimum is attained. The specifically stated question becomes
\[
 \forall C\in\mathcal A_{\mathrm{ss}},\qquad H_{\mathrm{sing}}(C)>-3.
\]
This pointwise strict inequality does not itself assert \(h_{\mathrm{ss}}>-3\). The source writes \(H(C)\); identifying that invariant with \(H_{\mathrm{sing}}(C)\), and adopting the arrangement classes above, remain unresolved source qualifications.
\subsection{Short English statement}
Determine the greatest possible negativity of the H-constant for rational line arrangements, complex line arrangements, and complex supersolvable line arrangements. In particular, decide whether every complex supersolvable arrangement satisfies \(H(C)>-3\). The displayed mathematical formulation is conditional because the supplied record does not specify the H-constant or the full arrangement conventions.
\subsection{Sources}
\begin{itemize}
\item Supplied local UnsolvedMath registry record, permanent ID \(30004604\), title \emph{Extremal H-Constants of Line Arrangements}, exact source problem code \texttt{catalogue entry}, in the Oberwolfach Reports Open Problems collection. The proposed year is recorded as 2021; no individual proposer is named.
\item Supplied report citation: \emph{Logarithmic Vector Fields and Freeness of Divisors and Arrangements: New perspectives and applications} (2022), Oberwolfach Reports, pp.~225--300, DOI 10.4171/owr/2021/5. \url{https://doi.org/10.4171/owr/2021/5}. The supplied literature list also labels this reference \emph{Oberwolfach Report 5/2021}.
\item Supplied literature reference: Urzua, \emph{Proyecciones} (2022). \url{https://www.scielo.cl/pdf/proy/v41n2/0716-0917-proy-41-02-517.pdf}. No fuller bibliographic identification is supplied.
\item Literature assessment dated exactly \texttt{2026-08-22}: the supplied curation labels the problem open and describes extremal-value and classification questions for rational and complex line arrangements, including the supersolvable subcase, as remaining open. This is supplied dated curation derived from the OWR record, not an independently verified literature assessment. The record expressly requests independent verification of time-sensitive status and cited claims.
\item The supplied statement status is \texttt{corrected\_verified}, accompanied by a report of manual semantic, typographic, and LaTeX review. This records the supplied editorial assessment; it does not independently establish the missing definitions or mathematical correctness of this draft.
\end{itemize}
\end{document}
